\subsection{A common finite-to-fluid transfer}
\label{sec:common-finite-transfer}

The fluid theorem deliberately fixed a known distribution and made
proportional revelation exact.  We now show that learning and private
shuffling reproduce both features up to $o(n^2)$ cost on bounded finite
inputs.  Doing this once also exposes the precise common structure of
obligatory testing, revealing optimization, and blind execution.

Fix $L<\infty$ and a known function
\[
 r:[0,L]\longrightarrow[0,L]
\]
with $|r(p)-r(p')|\le|p-p'|$.  A job first touched in the \emph{direct mode}
runs nonpreemptively for $r(p)$ and completes.  The choice of this mode must
be made before $p$ is known; revealing $p$ after completion may be granted
to the algorithm for free.  Alternatively, the job may be tested for one
unit, after which its processing operation of length $p$ can be scheduled.

Let $Q$ be either $[0,1]$, when both first-touch modes are available, or
$\{1\}$, when every job must first be tested.  For a distribution $D$, put
\begin{equation}
 v_r(D)\defeq\int r(p)\,dD(p).
 \label{eq:common-direct-mean}
\end{equation}
Predictably chosen direct jobs have asymptotic mean work $v_r(D)$ per job.
We therefore define the common fluid value
\begin{equation}
 \Phi_{r,Q}(D)\defeq
 \min_{q\in Q}\Area\bigl(\mathcal C^{v_r(D)}_{D,q}\bigr).
 \label{eq:common-instance-benchmark}
\end{equation}
When $Q=\{1\}$, the alternative block has zero mass and $r$ is immaterial.

The three specializations used later are
\begin{equation}
 \begin{array}{c|c|c}
  \text{model}&r(p)&Q\\
  \hline
  \mathsf{OT}&0&\{1\}\\
  \mathsf{RO}(u)&u&[0,1]\\
  \mathsf{BE}&p&[0,1].
 \end{array}
 \label{eq:common-instance-specializations}
\end{equation}
Thus the distinction between the last two models is entirely in the direct
work: it is fixed in revealing optimization and sampled from the input in
blind execution.

\begin{theorem}[Common bounded instance-optimality transfer]
\label{thm:common-fluid-instance-transfer}
Fix $L$, $r$, and $Q$ as above.  There are randomized nonanticipating
algorithms $\mathcal L_n$, depending only on $n,L,r,Q$, and a deterministic
sequence
\begin{equation}
 \eps_L(n)
 =O_L\!\left(n^{-1/6}\sqrt{\ln(n+2)}\right)
 \longrightarrow0
 \label{eq:common-transfer-rate}
\end{equation}
such that every vector $p\in[0,L]^n$ satisfies
\begin{equation}
 \EE\ALG_{\mathcal L_n}(p)
 \le n^2\Phi_{r,Q}(D[p])+\eps_L(n)n^2.
 \label{eq:common-transfer-upper}
\end{equation}
Conversely, every randomized nonanticipating algorithm $\mathcal A$ whose
allowed first-touch modes are described by $Q$, even if it is told the
multiset of entries of $p$, satisfies
\begin{equation}
 \EE\ALG_{\mathcal A}(p)
 \ge n^2\Phi_{r,Q}(D[p])-\eps_L(n)n^2.
 \label{eq:common-transfer-lower}
\end{equation}
Both expectations include the private shuffle and the algorithms' remaining
randomness.
\end{theorem}

The proof is split into a lower transfer and a learning upper transfer.  The
first maps every finite adaptive run to the envelope from
\cref{eq:common-fluid-envelope}.  The second estimates the empirical
distribution and implements the finite block schedule
$\textsc{FluidSPT}$.  The theorem then follows immediately by using the same
fluid value between them.

\subsection{From an adaptive run to the fluid envelope}
\label{sec:common-lower-transfer}

The central fact is that adaptivity cannot bias the hidden value at a first
touch.  After fixing an algorithm's private seed, list jobs in the order in
which their labels are first touched.  Their values form a uniform
permutation $X_1,\ldots,X_n$ of the input.  The indicator $A_k$ that the
$k$th job is tested is determined before $X_k$ is exposed.  Hence $A_k$ and
$1-A_k$ are predictable selectors in
\cref{lem:predictable-permutation-sampling}.

Keep zero as a separate class, partition $(0,L]$ into
\begin{equation}
 K_n=\lfloor n^{1/6}\rfloor
 \label{eq:common-transfer-grid-size}
\end{equation}
cells, and round positive values upward to their cell endpoints.  Write
$D^+$ for the resulting empirical distribution and $h_n=L/K_n$ for the
mesh.  Take $\delta_n=n^{-1/6}$ and stop the first-touch order while
$\delta_n n+O(1)$ labels remain untouched.

Apply the predictable-permutation estimate simultaneously to the vector of
cell indicators selected by $A_k$ and to the scalar $r(X_k)$ selected by
$1-A_k$.  With probability at least $1-n^{-2}$, every earlier prefix obeys
\begin{align}
 \bigl|\#\{\text{tested jobs in cell }j\}
      -d_j\#\{\text{tests}\}\bigr|&\le\Gamma_n,
 \label{eq:common-transfer-cell-discrepancy}\\
 \bigl|\text{direct work}
      -v_r(D[p])\#\{\text{direct jobs}\}\bigr|&\le\Gamma_n,
 \label{eq:common-transfer-work-discrepancy}
\end{align}
for all cells and all such prefixes, where
\begin{equation}
 \frac{(K_n+1)\Gamma_n}{n}
 =O\!\left(n^{-1/6}\sqrt{\ln(n+2)}\right).
 \label{eq:common-transfer-discrepancy-rate}
\end{equation}
For $Q=\{1\}$ the second estimate is vacuous.

Edit the last $\delta_n n+O(1)$ first touches to one fixed legal mode: test
them when $Q=\{1\}$ and complete them directly when $Q=[0,1]$.  Delete any
operations superseded by this edit.  At most $\delta_n n+O(1)$ jobs and a
constant number of operations per job change.  Since every legal job uses at
most $1+L$ work, operation-by-operation remaining-job accounting gives
\begin{equation}
 C_{\rm original}\ge C_{\rm edited}
   -O_L(\delta_n n^2+n).
 \label{eq:common-transfer-suffix-edit}
\end{equation}
Let $q$ be the final tested fraction of the edited run; then $q\in Q$.

\begin{claim}[Repairing a finite prefix]
\label{clm:common-envelope-repair}
On the event
\cref{eq:common-transfer-cell-discrepancy,eq:common-transfer-work-discrepancy},
there is
\begin{equation}
 \zeta_n=O_L\!\left(
   \frac{(K_n+1)\Gamma_n}{n}+h_n+\delta_n\right)
 \label{eq:common-transfer-zeta}
\end{equation}
such that, simultaneously throughout the edited run,
\begin{equation}
 \text{completed fraction after normalized work }x
 \le \mathcal C^{v_r(D^+)}_{D^+,q}(x+\zeta_n)+\zeta_n.
 \label{eq:common-transfer-envelope-repair}
\end{equation}
\end{claim}

\begin{proof}
At an operation boundary, let $t$ be tested mass, $b$ directly completed
mass, and $c_j$ the processed tested mass in positive cell $j$.  The cell
estimate gives
\[
 c_j\le d_jt+\Gamma_n/n,
\]
and gives the analogous bound for tested zeros.  Replace each $c_j$ by
$(c_j-\Gamma_n/n)^+$.  The resulting class capacities are exact and the
removed completed mass is at most $(K_n+1)\Gamma_n/n$.

The work estimate replaces the actual direct work by
$v_r(D[p])b$ with error at most $\Gamma_n/n$.  Upward rounding changes tested
processing work by at most $h_n$ and, because $r$ is one-Lipschitz, changes
the direct mean by at most $h_n$.  The edited suffix has mass
$O(\delta_n)$.  After these horizontal repairs, the variables $t,b,c_j$
satisfy every constraint in \cref{eq:common-fluid-envelope} with
$v=v_r(D^+)$.  This proves the displayed inequality with
\cref{eq:common-transfer-zeta}.
\end{proof}

We shall also use the following elementary stability observation.  It avoids
any claim that the optimizing threshold itself changes continuously.

\begin{lemma}[Zero-preserving rounding stability]
\label{lem:common-transfer-rounding-stability}
Suppose $D,D'$ admit a coupling that preserves zero and satisfies
$|P-P'|\le h$ almost surely.  Then
\begin{equation}
 |\Phi_{r,Q}(D)-\Phi_{r,Q}(D')|\le2h.
 \label{eq:common-transfer-rounding-stability}
\end{equation}
\end{lemma}

\begin{proof}
Use the same tested fraction, selectors, and block order for coupled mass
elements.  At every prefix, at most one unit of tested mass receives
processing, so its work changes by at most $h$.  Direct work per unit mass
also changes by at most $h$ because $r$ is one-Lipschitz.  Thus either
completion curve lies within a horizontal $2h$ shift of the other.  Their
areas differ by at most $2h$.  Exchange $D,D'$ and minimize over schedules.
\end{proof}

\begin{claim}[Common announced lower bound]
\label{clm:common-announced-lower}
Every announced randomized algorithm in the common finite model satisfies
\begin{equation}
 \frac{\EE\ALG_{\mathcal A}(p)}{n^2}
 \ge\Phi_{r,Q}(D[p])
 -O_L\!\left(n^{-1/6}\sqrt{\ln(n+2)}\right).
 \label{eq:common-announced-lower}
\end{equation}
\end{claim}

\begin{proof}
The envelope reaches one after at most $1+L$ normalized work.  Integrating
\cref{eq:common-transfer-envelope-repair} therefore loses only
$O_L(\zeta_n)$ from its area.  By
\cref{lem:test-module-envelope,eq:common-instance-benchmark}, this area is at
least $\Phi_{r,Q}(D^+)$.  The sampling event fails with probability at most
$n^{-2}$, while every fluid value is at most $(1+L)/2$.  Finally apply the
suffix edit and \cref{lem:common-transfer-rounding-stability}.  Equations
\eqref{eq:common-transfer-discrepancy-rate} and
\eqref{eq:common-transfer-zeta} give the stated rate.  The argument is
uniform after fixing the algorithm seed, so averaging proves the randomized
claim.
\end{proof}

\subsection{Learning and implementing the fluid optimum}
\label{sec:common-learning-transfer}

For the upper transfer, a \emph{grid template} records a tested fraction, a
selector on the grid classes, and a feasible order of the resulting testing,
deferred, and direct blocks.  For a grid distribution $D$ and template
$\pi$, let $F_{r,D}(\pi)$ be its fluid completion cost.  Theorem
\ref{lem:test-module-envelope} says that
\begin{equation}
 \min_\pi F_{r,D}(\pi)=\Phi_{r,Q}(D).
 \label{eq:common-template-optimum}
\end{equation}

Two finite facts make this template learnable.  We state them separately so
that the model-specific corollaries can reuse them.

\begin{claim}[Finite implementation of a fixed template]
\label{clm:common-finite-template-kernel}
For every fixed grid template $\pi$, selecting the prescribed number of test
labels uniformly and executing its blocks gives
\begin{equation}
 \EE C_n(\pi)=n^2F_{r,D[p]}(\pi)+O_L(n),
 \label{eq:common-finite-template-kernel}
\end{equation}
uniformly over the grid, its frequencies, and $\pi$.
\end{claim}

\begin{proof}
Expand total completion time into one-job and unordered-pair contributions.
For two distinct jobs, the probabilities that zero, one, or both are tested
differ from the corresponding products of the tested fraction by $O(1/n)$.
Every operation and pair contribution is bounded in terms of $L$.
Summing the diagonal terms and these finite-population corrections gives
$O_L(n)$; the remaining terms are exactly the fluid block area.
\end{proof}

\begin{claim}[Uniform stability of fixed templates]
\label{clm:common-template-stability}
There is $C_L<\infty$ such that any two distributions $D,E$ on the same
grid and every common template $\pi$ satisfy
\begin{equation}
 |F_{r,D}(\pi)-F_{r,E}(\pi)|
 \le C_L\|D-E\|_1.
 \label{eq:common-template-stability}
\end{equation}
The same inequality holds after minimizing over templates.
\end{claim}

\begin{proof}
For a fixed template, the block-area expansion is a finite sum of single and
double integrals whose kernels are bounded by a constant depending only on
$L$.  In a double integral, replace the first distribution and then the
second.  This gives the displayed total-variation bound uniformly over the
number of grid classes and over the template.  Applying the bound first to a
minimizer for $D$ and then to one for $E$ proves the assertion for minima.
\end{proof}

The learner uses the grid from
\cref{eq:common-transfer-grid-size}, privately permutes the labels, and tests
the first
\begin{equation}
 k_n=\lceil n^{1/2}\rceil
 \label{eq:common-transfer-pilot-size}
\end{equation}
jobs.  It keeps positive sample jobs pending and forms their rounded
histogram $\widehat D_n$.  It chooses a template minimizing
$F_{r,\widehat D_n}$ and applies it to the unused labels, then finishes the
positive sample jobs.

Let $D_n$ and $D_n^{\rm rem}$ be the full and remaining rounded histograms.
The sampling-without-replacement estimate
\cref{eq:without-replacement-histogram} and deletion of the pilot give
\begin{equation}
 \EE\|\widehat D_n-D_n\|_1
 \le\sqrt{\frac{K_n+1}{k_n}},
 \qquad
 \|D_n^{\rm rem}-D_n\|_1
 \le\frac{2k_n}{n-k_n}.
 \label{eq:common-transfer-histogram}
\end{equation}

\begin{claim}[Common learning upper bound]
\label{clm:common-learning-upper}
The learner above satisfies
\begin{equation}
 \frac{\EE\ALG_{\mathcal L_n}(p)}{n^2}
 \le\Phi_{r,Q}(D[p])+O_L(n^{-1/6})
 \label{eq:common-learning-upper}
\end{equation}
uniformly over $p\in[0,L]^n$.
\end{claim}

\begin{proof}
Condition on the pilot.  The unused labels remain uniformly permuted, so
\cref{clm:common-finite-template-kernel} applies to the selected template.
The pilot and its final positive tail use $O_L(k_n)$ work and therefore add
$O_L(nk_n)$ to total completion time.

The fixed-template kernel is applied to the upward-rounded values.  Coupling
the physical and rounded executions with the same decisions changes tested
processing work by at most $h_n$ per unit mass and direct work by at most
$h_n$ per unit mass, because $r$ is one-Lipschitz.  Remaining-job accounting
therefore changes the normalized template cost by at most $O(h_n)$.

Apply \cref{clm:common-template-stability} first to the learned template,
then to an optimal template for $D_n^{\rm rem}$, and finally to an optimal
template for $D_n$.  This bounds the conditional fluid cost by
\[
 \Phi_{r,Q}(D_n)
 +O_L\bigl(
   \|\widehat D_n-D_n\|_1+
   \|D_n^{\rm rem}-D_n\|_1\bigr).
\]
Average and use \cref{eq:common-transfer-histogram}.  The three remaining
errors are
\[
 h_n=O_L(n^{-1/6}),\qquad
 \sqrt{(K_n+1)/k_n}=O(n^{-1/6}),\qquad
 k_n/n=O(n^{-1/2}).
\]
The rounding term is controlled by
\cref{lem:common-transfer-rounding-stability}, proving the claim.
\end{proof}

\begin{proof}[Proof of \cref{thm:common-fluid-instance-transfer}]
Claim~\ref{clm:common-learning-upper} gives
\cref{eq:common-transfer-upper}.  Claim~\ref{clm:common-announced-lower}
gives \cref{eq:common-transfer-lower}, since announcing the multiset only
strengthens the competing algorithm.  Taking the larger of their vanishing
errors gives \cref{eq:common-transfer-rate}.
\end{proof}
